\documentclass[11pt]{article}
\usepackage[a4paper,margin=2cm]{geometry}
\usepackage[T1]{fontenc}
\usepackage[utf8]{inputenc}
\usepackage{lmodern}
\usepackage{enumitem}
\usepackage{parskip}
\setlength{\parindent}{0pt}
\setlength{\fboxsep}{8pt}
\setlength{\fboxrule}{0.5pt}
\renewcommand{\arraystretch}{1.2}

\begin{document}

\section*{Human Evaluation Sheet}
% Write your own intro above if needed.


% item=1 | prompt_id=78 | iteration=3 | compare_mode=base_vs_retrieved | judge_winner=base | A_source=Retrieved | B_source=Base
\subsubsection*{Item 1}
\textbf{Prompt ID:} 78 \hfill \textbf{Iteration:} 3

\textbf{Task Prompt:} Detect and caption 'examples/family\_dinner.jpg' and narrate it.

\vspace{0.75em}
\noindent
\begin{minipage}[t]{0.48\textwidth}
\textbf{Explanation A}

\vspace{0.5em}
\fbox{
\parbox[t][0.30\textheight][t]{0.95\linewidth}{To achieve this, I used a two-stage approach. First, I utilized an image classification model to identify the content of the image. The model used was 'microsoft/beit-base-patch16-224-pt22k-ft22k', which provided predictions such as 'King's\_Counsel', 'dinner\_table', 'breakfast\_table', 'kitchen\_table', and 'buffet' with confidence scores. However, since the task was to detect and caption the image, I relied more on the image-to-text model for the final caption. The image-to-text model 'nlpconnect/vit-gpt2-image-captioning' generated the caption "a family gathers around a table for dinner" directly from the image. This caption aligns well with the context of a family dinner scene. The generated audio of the caption can be found at "[PATH]".}
}
\end{minipage}
\hfill
\begin{minipage}[t]{0.48\textwidth}
\textbf{Explanation B}

\vspace{0.5em}
\fbox{
\parbox[t][0.30\textheight][t]{0.95\linewidth}{To achieve this, I used a two-stage approach. First, I utilized an image classification model to identify the content of the image. The model used was 'microsoft/beit-base-patch16-224-pt22k-ft22k', which provided predictions such as 'King's\_Counsel', 'dinner\_table', 'breakfast\_table', 'kitchen\_table', and 'buffet' with confidence scores. However, since the task was to detect and caption the image, I relied more on the image-to-text model for the final caption. The image-to-text model 'nlpconnect/vit-gpt2-image-captioning' generated the caption "a family gathers around a table for dinner" directly from the image. This caption aligns well with the context of a family dinner scene. The generated audio of the caption can be found at "/audios/9e96.wav".}
}
\end{minipage}

\vspace{1em}

\textbf{Evaluation}
\begin{itemize}
    \item \textbf{Better explanation:} \hspace{1em} $\square$ A \hspace{1em} $\square$ B \hspace{1em} $\square$ Equal
    \item \textbf{Grade A:} \underline{\hspace{2cm}} \hspace{1em} \textbf{Grade B:} \underline{\hspace{2cm}}
    \item \textbf{Reason:} \\[4.5em]
\end{itemize}

\vspace{1em}
\hrule
\vspace{1em}


% item=2 | prompt_id=75 | iteration=6 | compare_mode=base_vs_combined | judge_winner=combined | A_source=Combined | B_source=Base
\subsubsection*{Item 2}
\textbf{Prompt ID:} 75 \hfill \textbf{Iteration:} 6

\textbf{Task Prompt:} Merge captions, detections, and segmentation to illustrate 'examples/beach\_party.jpg'.

\vspace{0.75em}
\noindent
\begin{minipage}[t]{0.48\textwidth}
\textbf{Explanation A}

\vspace{0.5em}
\fbox{
\parbox[t][0.30\textheight][t]{0.95\linewidth}{Given the inconsistency in the image-to-text task's output and the expected content of the image, it seems there might have been a misunderstanding or misclassification in the process. The object detection task successfully identified several objects relevant to a beach party scenario, such as 17 people, 5 boats, and 4 umbrellas, with high confidence scores ranging from 0.9037 to 0.9972 for people, 0.9556 to 0.9920 for boats, and 0.9450 to 0.9765 for umbrellas. However, the final caption generated does not align with these detections, suggesting a potential error in the image-to-text task or in the integration of the tasks' outputs. The models used in the process included microsoft/beit-base-patch16-224-pt22k-ft22k for image classification, facebook/detr-resnet-101 for object detection, facebook/maskformer-swin-base-coco for image segmentation, and nlpconnect/vit-gpt2-image-captioning for the image-to-text task. The object detection task provided a generated image with bounding boxes around the detected objects, and the image segmentation task provided a segmented image of the scene. Given the inconsistency in the results, particularly with the final caption not matching the expected beach party theme, it's challenging to provide a coherent merged output that accurately represents the image.}
}
\end{minipage}
\hfill
\begin{minipage}[t]{0.48\textwidth}
\textbf{Explanation B}

\vspace{0.5em}
\fbox{
\parbox[t][0.30\textheight][t]{0.95\linewidth}{Given the inconsistency in the image-to-text task's output and the expected content of the image, it seems there might have been a misunderstanding or misclassification in the process. The object detection task successfully identified several objects relevant to a beach party scenario, such as people, umbrellas, and boats, with high confidence scores. However, the final caption generated does not align with these detections, suggesting a potential error in the image-to-text task or in the integration of the tasks' outputs. 
The models used in the process included "microsoft/beit-base-patch16-224-pt22k-ft22k" for image classification, "facebook/detr-resnet-101" for object detection, "facebook/maskformer-swin-base-coco" for image segmentation, and "nlpconnect/vit-gpt2-image-captioning" for the image-to-text task. 
The object detection task provided a generated image located at "/images/4fcd.jpg", and the image segmentation task provided a segmented image at "/images/6b44.jpg". 
Given the inconsistency in the results, particularly with the final caption not matching the expected beach party theme, it's challenging to provide a coherent merged output that accurately represents the image 'examples/beach\_party.jpg'.}
}
\end{minipage}

\vspace{1em}

\textbf{Evaluation}
\begin{itemize}
    \item \textbf{Better explanation:} \hspace{1em} $\square$ A \hspace{1em} $\square$ B \hspace{1em} $\square$ Equal
    \item \textbf{Grade A:} \underline{\hspace{2cm}} \hspace{1em} \textbf{Grade B:} \underline{\hspace{2cm}}
    \item \textbf{Reason:} \\[4.5em]
\end{itemize}

\vspace{1em}
\hrule
\vspace{1em}


% item=3 | prompt_id=11 | iteration=2 | compare_mode=base_vs_retrieved | judge_winner=retrieved | A_source=Base | B_source=Retrieved
\subsubsection*{Item 3}
\textbf{Prompt ID:} 11 \hfill \textbf{Iteration:} 2

\textbf{Task Prompt:} Extract total amount from 'examples/invoice\_sample.jpg'.

\vspace{0.75em}
\noindent
\begin{minipage}[t]{0.48\textwidth}
\textbf{Explanation A}

\vspace{0.5em}
\fbox{
\parbox[t][0.30\textheight][t]{0.95\linewidth}{To fulfill your request, I attempted to use a series of models. First, I used an image classification model to analyze the image 'examples/invoice\_sample.jpg'. However, this model did not provide any relevant information for extracting the total amount. Then, I tried to use an optical character recognition (OCR) model to extract text from the image, but unfortunately, this task was not available. Without the OCR results, I couldn't proceed with the document question answering model to find the total amount.}
}
\end{minipage}
\hfill
\begin{minipage}[t]{0.48\textwidth}
\textbf{Explanation B}

\vspace{0.5em}
\fbox{
\parbox[t][0.30\textheight][t]{0.95\linewidth}{To fulfill your request, I attempted to use a series of models. First, I used an image classification model to analyze the image '[PATH]'. However, this model did not provide any relevant information for extracting the total amount. Then, I tried to use an optical character recognition (OCR) model to extract text from the image, but unfortunately, this task was not available. Without the OCR results, I couldn't proceed with the document question answering model to find the total amount. To proceed, consider checking the availability of the OCR task or exploring alternative methods for text extraction from the image.}
}
\end{minipage}

\vspace{1em}

\textbf{Evaluation}
\begin{itemize}
    \item \textbf{Better explanation:} \hspace{1em} $\square$ A \hspace{1em} $\square$ B \hspace{1em} $\square$ Equal
    \item \textbf{Grade A:} \underline{\hspace{2cm}} \hspace{1em} \textbf{Grade B:} \underline{\hspace{2cm}}
    \item \textbf{Reason:} \\[4.5em]
\end{itemize}

\vspace{1em}
\hrule
\vspace{1em}


% item=4 | prompt_id=39 | iteration=7 | compare_mode=base_vs_combined | judge_winner=combined | A_source=Base | B_source=Combined
\subsubsection*{Item 4}
\textbf{Prompt ID:} 39 \hfill \textbf{Iteration:} 7

\textbf{Task Prompt:} Identify the pose in 'examples/runner\_pose.jpg' and create an image of 'a person jumping'.

\vspace{0.75em}
\noindent
\begin{minipage}[t]{0.48\textwidth}
\textbf{Explanation A}

\vspace{0.5em}
\fbox{
\parbox[t][0.30\textheight][t]{0.95\linewidth}{I used the openpose-control model to identify the pose in the given image, but unfortunately, the model was not loaded and returned an error. However, I was still able to generate an image of a person jumping using the openpose-text-to-image model, specifically the lllyasviel/sd-controlnet-openpose model, with the text prompt "a person jumping" and the image "examples/runner\_pose.jpg". The generated image is located at /images/becb.png.}
}
\end{minipage}
\hfill
\begin{minipage}[t]{0.48\textwidth}
\textbf{Explanation B}

\vspace{0.5em}
\fbox{
\parbox[t][0.30\textheight][t]{0.95\linewidth}{I used the openpose-control model to identify the pose in the given image, but unfortunately, the model was not loaded and returned an error. However, I was still able to generate an image of a person jumping using the lllyasviel/sd-controlnet-openpose model, which is a ControlNet model that can generate images based on text prompts and an input image. I chose to use the ControlNet model because it is well-suited for tasks that require generating images based on text prompts and input images. In this case, I used the text prompt 'a person jumping' and the image 'examples/runner\_pose.jpg'. The generated image is located at /images/becb.png. Please note that the accuracy of the pose identification and image generation may vary depending on the quality of the input image and the complexity of the task.}
}
\end{minipage}

\vspace{1em}

\textbf{Evaluation}
\begin{itemize}
    \item \textbf{Better explanation:} \hspace{1em} $\square$ A \hspace{1em} $\square$ B \hspace{1em} $\square$ Equal
    \item \textbf{Grade A:} \underline{\hspace{2cm}} \hspace{1em} \textbf{Grade B:} \underline{\hspace{2cm}}
    \item \textbf{Reason:} \\[4.5em]
\end{itemize}

\vspace{1em}
\hrule
\vspace{1em}


% item=5 | prompt_id=9 | iteration=1 | compare_mode=base_vs_retrieved | judge_winner=base | A_source=Base | B_source=Retrieved
\subsubsection*{Item 5}
\textbf{Prompt ID:} 9 \hfill \textbf{Iteration:} 1

\textbf{Task Prompt:} Answer 'What is the person holding?' in the image 'examples/club.jpg'.

\vspace{0.75em}
\noindent
\begin{minipage}[t]{0.48\textwidth}
\textbf{Explanation A}

\vspace{0.5em}
\fbox{
\parbox[t][0.30\textheight][t]{0.95\linewidth}{To answer your question, I used a two-stage approach. First, I utilized the 'nlpconnect/vit-gpt2-image-captioning' model to generate a caption for the image. The model generated the text: "a woman holding a glass of beer next to a group of people". Then, I used the 'dandelin/vilt-b32-finetuned-vqa' model to specifically answer the question "What is the person holding?" based on the image and the question. The model's answer was "beer", which aligns with the generated caption. Since both models provided consistent results, I concluded that the person is holding a glass of beer. The image used for the inference is located at 'public/examples/club.jpg'.}
}
\end{minipage}
\hfill
\begin{minipage}[t]{0.48\textwidth}
\textbf{Explanation B}

\vspace{0.5em}
\fbox{
\parbox[t][0.30\textheight][t]{0.95\linewidth}{To answer your question, I used a two-stage approach. First, I utilized the 'nlpconnect/vit-gpt2-image-captioning' model to generate a caption for the image. The model generated the text: "a woman holding a glass of beer next to a group of people". Then, I used the 'dandelin/vilt-b32-finetuned-vqa' model to specifically answer the question "What is the person holding?" based on the image and the question. The model's answer was "beer", which aligns with the generated caption. Since both models provided consistent results, I concluded that the person is holding a glass of beer. The image used for the inference is located at 'public/examples/club.jpg'.}
}
\end{minipage}

\vspace{1em}

\textbf{Evaluation}
\begin{itemize}
    \item \textbf{Better explanation:} \hspace{1em} $\square$ A \hspace{1em} $\square$ B \hspace{1em} $\square$ Equal
    \item \textbf{Grade A:} \underline{\hspace{2cm}} \hspace{1em} \textbf{Grade B:} \underline{\hspace{2cm}}
    \item \textbf{Reason:} \\[4.5em]
\end{itemize}

\vspace{1em}
\hrule
\vspace{1em}


% item=6 | prompt_id=69 | iteration=7 | compare_mode=base_vs_combined | judge_winner=combined | A_source=Combined | B_source=Base
\subsubsection*{Item 6}
\textbf{Prompt ID:} 69 \hfill \textbf{Iteration:} 7

\textbf{Task Prompt:} Identify objects and carry out segmentation on 'examples/warehouse.jpg'.

\vspace{0.75em}
\noindent
\begin{minipage}[t]{0.48\textwidth}
\textbf{Explanation A}

\vspace{0.5em}
\fbox{
\parbox[t][0.30\textheight][t]{0.95\linewidth}{To accomplish this task, I used a two-stage approach. First, I utilized the facebook/detr-resnet-101 model for object detection. This model detected 5 people in the image with confidence scores ranging from 0.9057 to 0.9978. The detection results are as follows: Person 1: score = 0.9308, box = (128, 452, 139, 468), Person 2: score = 0.9978, box = (198, 491, 219, 546), Person 3: score = 0.9807, box = (54, 446, 66, 478), Person 4: score = 0.9057, box = (337, 411, 347, 434), Person 5: score = 0.9673, box = (364, 419, 374, 441). The detection results are stored in the file /inference-results/examples/warehouse-detection.json. Then, I used the facebook/maskformer-swin-base-coco model for image segmentation, which generated a segmented image based on the detection results. The segmented image is saved at /inference-results/examples/warehouse-segmentation.png. I chose these models because they are state-of-the-art models for object detection and image segmentation tasks, respectively. The confidence scores and detection results suggest that the models performed well on this task. However, it's worth noting that the accuracy of the results may vary depending on the quality of the input image and the complexity of the scene.}
}
\end{minipage}
\hfill
\begin{minipage}[t]{0.48\textwidth}
\textbf{Explanation B}

\vspace{0.5em}
\fbox{
\parbox[t][0.30\textheight][t]{0.95\linewidth}{To accomplish this task, I used a two-stage approach. First, I utilized the facebook/detr-resnet-101 model for object detection. This model detected 5 people in the image with confidence scores ranging from 0.9057 to 0.9978. The detection results are stored in the file /inference-results/examples/warehouse-detection.json. Then, I used the facebook/maskformer-swin-base-coco model for image segmentation, which generated a segmented image based on the detection results. The segmented image is saved at /inference-results/examples/warehouse-segmentation.png. I chose these models because they are state-of-the-art models for object detection and image segmentation tasks, respectively. The confidence scores and detection results suggest that the models performed well on this task. However, it's worth noting that the accuracy of the results may vary depending on the quality of the input image and the complexity of the scene.}
}
\end{minipage}

\vspace{1em}

\textbf{Evaluation}
\begin{itemize}
    \item \textbf{Better explanation:} \hspace{1em} $\square$ A \hspace{1em} $\square$ B \hspace{1em} $\square$ Equal
    \item \textbf{Grade A:} \underline{\hspace{2cm}} \hspace{1em} \textbf{Grade B:} \underline{\hspace{2cm}}
    \item \textbf{Reason:} \\[4.5em]
\end{itemize}

\vspace{1em}
\hrule
\vspace{1em}


% item=7 | prompt_id=52 | iteration=5 | compare_mode=base_vs_retrieved | judge_winner=retrieved | A_source=Base | B_source=Retrieved
\subsubsection*{Item 7}
\textbf{Prompt ID:} 52 \hfill \textbf{Iteration:} 5

\textbf{Task Prompt:} Translate '人工智能改变世界' into English.

\vspace{0.75em}
\noindent
\begin{minipage}[t]{0.48\textwidth}
\textbf{Explanation A}

\vspace{0.5em}
\fbox{
\parbox[t][0.30\textheight][t]{0.95\linewidth}{The translation of '人工智能改变世界' into English is:

"Artificial Intelligence Changes the World"}
}
\end{minipage}
\hfill
\begin{minipage}[t]{0.48\textwidth}
\textbf{Explanation B}

\vspace{0.5em}
\fbox{
\parbox[t][0.30\textheight][t]{0.95\linewidth}{The translation of '人工智能改变世界' into English is 'Artificial Intelligence Changes the World'.}
}
\end{minipage}

\vspace{1em}

\textbf{Evaluation}
\begin{itemize}
    \item \textbf{Better explanation:} \hspace{1em} $\square$ A \hspace{1em} $\square$ B \hspace{1em} $\square$ Equal
    \item \textbf{Grade A:} \underline{\hspace{2cm}} \hspace{1em} \textbf{Grade B:} \underline{\hspace{2cm}}
    \item \textbf{Reason:} \\[4.5em]
\end{itemize}

\vspace{1em}
\hrule
\vspace{1em}


% item=8 | prompt_id=25 | iteration=6 | compare_mode=base_vs_retrieved | judge_winner=base | A_source=Base | B_source=Retrieved
\subsubsection*{Item 8}
\textbf{Prompt ID:} 25 \hfill \textbf{Iteration:} 6

\textbf{Task Prompt:} Categorize the items in 'examples/bridge.png'.

\vspace{0.75em}
\noindent
\begin{minipage}[t]{0.48\textwidth}
\textbf{Explanation A}

\vspace{0.5em}
\fbox{
\parbox[t][0.30\textheight][t]{0.95\linewidth}{To determine the categorization of the items in 'examples/bridge.png', I utilized a single model, 'microsoft/beit-base-patch16-224-pt22k-ft22k', for image classification. This model provided several predictions with their corresponding confidence scores. The top predictions included 'pied-a-terre' with a score of 0.1606227308511734, 'gate' with a score of 0.16047406196594238, and 'cantilever\_bridge' with a score of 0.11095928400754929. Given these results, it seems the image could be related to a bridge or structures associated with it, but the model's confidence in any single category is not extremely high. The image path used for this inference is 'public/examples/bridge.png'. Unfortunately, without more context or higher confidence scores, it's challenging to provide a more specific categorization.}
}
\end{minipage}
\hfill
\begin{minipage}[t]{0.48\textwidth}
\textbf{Explanation B}

\vspace{0.5em}
\fbox{
\parbox[t][0.30\textheight][t]{0.95\linewidth}{To determine the categorization of the items in 'examples/bridge.png', I utilized a single model, 'microsoft/beit-base-patch16-224-pt22k-ft22k', for image classification. This model provided several predictions with their corresponding confidence scores. The top predictions included 'pied-a-terre' with a score of 0.1606227308511734, 'gate' with a score of 0.16047406196594238, and 'cantilever\_bridge' with a score of 0.11095928400754929. Given these results, it seems the image could be related to a bridge or structures associated with it, but the model's confidence in any single category is not extremely high. The image path used for this inference is 'public/examples/bridge.png'. Unfortunately, without more context or higher confidence scores, it's challenging to provide a more specific categorization.}
}
\end{minipage}

\vspace{1em}

\textbf{Evaluation}
\begin{itemize}
    \item \textbf{Better explanation:} \hspace{1em} $\square$ A \hspace{1em} $\square$ B \hspace{1em} $\square$ Equal
    \item \textbf{Grade A:} \underline{\hspace{2cm}} \hspace{1em} \textbf{Grade B:} \underline{\hspace{2cm}}
    \item \textbf{Reason:} \\[4.5em]
\end{itemize}

\vspace{1em}
\hrule
\vspace{1em}


% item=9 | prompt_id=39 | iteration=3 | compare_mode=base_vs_retrieved | judge_winner=retrieved | A_source=Base | B_source=Retrieved
\subsubsection*{Item 9}
\textbf{Prompt ID:} 39 \hfill \textbf{Iteration:} 3

\textbf{Task Prompt:} Detect pose in 'examples/runner\_pose.jpg' and generate image 'a person jumping'.

\vspace{0.75em}
\noindent
\begin{minipage}[t]{0.48\textwidth}
\textbf{Explanation A}

\vspace{0.5em}
\fbox{
\parbox[t][0.30\textheight][t]{0.95\linewidth}{Given the limitations, I couldn't accurately detect the pose in the provided image. For the generation task, the model produced an image located at /images/ad50.jpg. 
The workflow started with an attempt to classify the image 'examples/runner\_pose.jpg' using the microsoft/beit-base-patch16-224-pt22k-ft22k model, which provided predictions but didn't directly contribute to pose detection or image generation. 
Then, it proceeded to the openpose-control task, which unfortunately failed due to the model not being loaded. 
Lastly, the runwayml/stable-diffusion-v1-5 model was used for generating the image based on the text 'a person jumping', resulting in the image /images/ad50.jpg.}
}
\end{minipage}
\hfill
\begin{minipage}[t]{0.48\textwidth}
\textbf{Explanation B}

\vspace{0.5em}
\fbox{
\parbox[t][0.30\textheight][t]{0.95\linewidth}{Given the provided image '[PATH]', I attempted to detect the pose using the openpose-control model, but unfortunately, it was not loaded, and I couldn't proceed with the pose detection. To resolve this issue, ensure the openpose-control model is properly loaded before attempting pose detection. Nevertheless, I was able to generate an image of 'a person jumping' using the runwayml/stable-diffusion-v1-5 model. The generated image is located at '[PATH]'. For further analysis, consider using the generated image as input for the pose detection task or explore other image classification models to gather more relevant information.}
}
\end{minipage}

\vspace{1em}

\textbf{Evaluation}
\begin{itemize}
    \item \textbf{Better explanation:} \hspace{1em} $\square$ A \hspace{1em} $\square$ B \hspace{1em} $\square$ Equal
    \item \textbf{Grade A:} \underline{\hspace{2cm}} \hspace{1em} \textbf{Grade B:} \underline{\hspace{2cm}}
    \item \textbf{Reason:} \\[4.5em]
\end{itemize}

\vspace{1em}
\hrule
\vspace{1em}


% item=10 | prompt_id=52 | iteration=2 | compare_mode=base_vs_retrieved | judge_winner=base | A_source=Retrieved | B_source=Base
\subsubsection*{Item 10}
\textbf{Prompt ID:} 52 \hfill \textbf{Iteration:} 2

\textbf{Task Prompt:} Translate '人工智能改变世界' into English.

\vspace{0.75em}
\noindent
\begin{minipage}[t]{0.48\textwidth}
\textbf{Explanation A}

\vspace{0.5em}
\fbox{
\parbox[t][0.30\textheight][t]{0.95\linewidth}{The translation of '人工智能改变世界' into English is 'Artificial Intelligence Changes the World'.}
}
\end{minipage}
\hfill
\begin{minipage}[t]{0.48\textwidth}
\textbf{Explanation B}

\vspace{0.5em}
\fbox{
\parbox[t][0.30\textheight][t]{0.95\linewidth}{The translation of '人工智能改变世界' into English is 'Artificial Intelligence Changes the World'.}
}
\end{minipage}

\vspace{1em}

\textbf{Evaluation}
\begin{itemize}
    \item \textbf{Better explanation:} \hspace{1em} $\square$ A \hspace{1em} $\square$ B \hspace{1em} $\square$ Equal
    \item \textbf{Grade A:} \underline{\hspace{2cm}} \hspace{1em} \textbf{Grade B:} \underline{\hspace{2cm}}
    \item \textbf{Reason:} \\[4.5em]
\end{itemize}

\vspace{1em}
\hrule
\vspace{1em}


% item=11 | prompt_id=38 | iteration=3 | compare_mode=base_vs_retrieved | judge_winner=retrieved | A_source=Retrieved | B_source=Base
\subsubsection*{Item 11}
\textbf{Prompt ID:} 38 \hfill \textbf{Iteration:} 3

\textbf{Task Prompt:} Generate a canny image for 'examples/bridge.png' then describe it.

\vspace{0.75em}
\noindent
\begin{minipage}[t]{0.48\textwidth}
\textbf{Explanation A}

\vspace{0.5em}
\fbox{
\parbox[t][0.30\textheight][t]{0.95\linewidth}{I used the lllyasviel/sd-controlnet-canny model to generate the Canny image, but it returned an error message indicating that the input should be an image. Then, I attempted to use the nlpconnect/vit-gpt2-image-captioning model to describe the generated image, but it also returned an error message. Given the errors, I couldn't obtain the confidence scores or output values for the models. To resolve the issue, I recommend checking the input format and ensuring it meets the model's requirements. If you'd like to try again or use a different model, please let me know, and I'll do my best to assist you.}
}
\end{minipage}
\hfill
\begin{minipage}[t]{0.48\textwidth}
\textbf{Explanation B}

\vspace{0.5em}
\fbox{
\parbox[t][0.30\textheight][t]{0.95\linewidth}{I used the lllyasviel/sd-controlnet-canny model to generate the Canny image, but it returned an error message indicating that the input should be an image. Then, I attempted to use the nlpconnect/vit-gpt2-image-captioning model to describe the generated image, but it also returned an error message. 
Given the errors, I couldn't obtain the confidence scores or output values for the models. The path to the generated image is not available since the image generation was unsuccessful. 
If you'd like to try again or use a different model, please let me know, and I'll do my best to assist you.}
}
\end{minipage}

\vspace{1em}

\textbf{Evaluation}
\begin{itemize}
    \item \textbf{Better explanation:} \hspace{1em} $\square$ A \hspace{1em} $\square$ B \hspace{1em} $\square$ Equal
    \item \textbf{Grade A:} \underline{\hspace{2cm}} \hspace{1em} \textbf{Grade B:} \underline{\hspace{2cm}}
    \item \textbf{Reason:} \\[4.5em]
\end{itemize}

\vspace{1em}
\hrule
\vspace{1em}


% item=12 | prompt_id=52 | iteration=3 | compare_mode=base_vs_retrieved | judge_winner=retrieved | A_source=Retrieved | B_source=Base
\subsubsection*{Item 12}
\textbf{Prompt ID:} 52 \hfill \textbf{Iteration:} 3

\textbf{Task Prompt:} Translate '人工智能改变世界' into English.

\vspace{0.75em}
\noindent
\begin{minipage}[t]{0.48\textwidth}
\textbf{Explanation A}

\vspace{0.5em}
\fbox{
\parbox[t][0.30\textheight][t]{0.95\linewidth}{The translation of '人工智能改变世界' into English is 'Artificial Intelligence Changes the World'.}
}
\end{minipage}
\hfill
\begin{minipage}[t]{0.48\textwidth}
\textbf{Explanation B}

\vspace{0.5em}
\fbox{
\parbox[t][0.30\textheight][t]{0.95\linewidth}{The translation of '人工智能改变世界' into English is:

"Artificial Intelligence Changes the World"}
}
\end{minipage}

\vspace{1em}

\textbf{Evaluation}
\begin{itemize}
    \item \textbf{Better explanation:} \hspace{1em} $\square$ A \hspace{1em} $\square$ B \hspace{1em} $\square$ Equal
    \item \textbf{Grade A:} \underline{\hspace{2cm}} \hspace{1em} \textbf{Grade B:} \underline{\hspace{2cm}}
    \item \textbf{Reason:} \\[4.5em]
\end{itemize}

\vspace{1em}
\hrule
\vspace{1em}


% item=13 | prompt_id=10 | iteration=5 | compare_mode=base_vs_retrieved | judge_winner=retrieved | A_source=Retrieved | B_source=Base
\subsubsection*{Item 13}
\textbf{Prompt ID:} 10 \hfill \textbf{Iteration:} 5

\textbf{Task Prompt:} Generate a video for the text 'A dog running through snow'.

\vspace{0.75em}
\noindent
\begin{minipage}[t]{0.48\textwidth}
\textbf{Explanation A}

\vspace{0.5em}
\fbox{
\parbox[t][0.30\textheight][t]{0.95\linewidth}{To create this video, I utilized a text-to-video model, specifically the damo-vilab/text-to-video-ms-1.7b model, which is designed to generate videos based on textual descriptions. This model was chosen because it is the most suitable option available for this task. The model processed the input text 'A dog running through snow' and produced a video that represents this scene. Although the model's confidence score is not available, the generated video can be reviewed for quality and accuracy, and if necessary, the model can be fine-tuned or different input text can be tried to achieve the desired outcome. The generated video can be found at the path [PATH] Please note that the video generation process is complex and the output may not always perfectly match the input text, but the model strives to create a visually coherent and relevant video.}
}
\end{minipage}
\hfill
\begin{minipage}[t]{0.48\textwidth}
\textbf{Explanation B}

\vspace{0.5em}
\fbox{
\parbox[t][0.30\textheight][t]{0.95\linewidth}{To generate the video, I used the damo-vilab/text-to-video-ms-1.7b model, which is a text-to-video model that can generate videos based on text descriptions. This model was chosen because it is the only available model for this task. The model took the text 'A dog running through snow' as input and generated a video that depicts a dog running through a snowy scene. The confidence score of the model is not available, but the generated video can be found at the path /videos/10e8.mp4.}
}
\end{minipage}

\vspace{1em}

\textbf{Evaluation}
\begin{itemize}
    \item \textbf{Better explanation:} \hspace{1em} $\square$ A \hspace{1em} $\square$ B \hspace{1em} $\square$ Equal
    \item \textbf{Grade A:} \underline{\hspace{2cm}} \hspace{1em} \textbf{Grade B:} \underline{\hspace{2cm}}
    \item \textbf{Reason:} \\[4.5em]
\end{itemize}

\vspace{1em}
\hrule
\vspace{1em}


% item=14 | prompt_id=25 | iteration=2 | compare_mode=base_vs_retrieved | judge_winner=base | A_source=Base | B_source=Retrieved
\subsubsection*{Item 14}
\textbf{Prompt ID:} 25 \hfill \textbf{Iteration:} 2

\textbf{Task Prompt:} Classify objects in 'examples/bridge.png'.

\vspace{0.75em}
\noindent
\begin{minipage}[t]{0.48\textwidth}
\textbf{Explanation A}

\vspace{0.5em}
\fbox{
\parbox[t][0.30\textheight][t]{0.95\linewidth}{To classify the objects in the given image, I utilized a single model, 'facebook/detr-resnet-101', as it was the only available option. The model detected 7 objects in total, with 2 cars, 4 boats, and 1 truck. The confidence scores for these detections ranged from 0.9177735447883606 to 0.9879690408706665, indicating a relatively high level of accuracy. However, it's essential to note that the model's performance may vary depending on the quality of the input image and the complexity of the scene. 
The generated image with bounding boxes can be found at '/images/a6c4.jpg'.}
}
\end{minipage}
\hfill
\begin{minipage}[t]{0.48\textwidth}
\textbf{Explanation B}

\vspace{0.5em}
\fbox{
\parbox[t][0.30\textheight][t]{0.95\linewidth}{To classify the objects in the given image, I used the facebook/detr-resnet-101 model, which is a popular and accurate model for object detection tasks. The model analyzed the image and provided the locations and labels of the detected objects. The confidence scores for each object were also provided, which helped me to determine the accuracy of the detections. Based on the results, I found that the image contains 2 cars, 1 truck, and 4 boats. The cars and truck are likely to be on the bridge, while the boats are on the water below. The model's confidence scores for these detections were high, ranging from 0.9177735447883606 to 0.9879690408706665, which suggests that the detections are accurate. The generated image with bounding boxes around the detected objects can be found at [PATH]}
}
\end{minipage}

\vspace{1em}

\textbf{Evaluation}
\begin{itemize}
    \item \textbf{Better explanation:} \hspace{1em} $\square$ A \hspace{1em} $\square$ B \hspace{1em} $\square$ Equal
    \item \textbf{Grade A:} \underline{\hspace{2cm}} \hspace{1em} \textbf{Grade B:} \underline{\hspace{2cm}}
    \item \textbf{Reason:} \\[4.5em]
\end{itemize}

\vspace{1em}
\hrule
\vspace{1em}


% item=15 | prompt_id=69 | iteration=2 | compare_mode=base_vs_combined | judge_winner=combined | A_source=Base | B_source=Combined
\subsubsection*{Item 15}
\textbf{Prompt ID:} 69 \hfill \textbf{Iteration:} 2

\textbf{Task Prompt:} Detect objects and perform segmentation on 'examples/warehouse.jpg'.

\vspace{0.75em}
\noindent
\begin{minipage}[t]{0.48\textwidth}
\textbf{Explanation A}

\vspace{0.5em}
\fbox{
\parbox[t][0.30\textheight][t]{0.95\linewidth}{To accomplish this task, I used a two-stage approach. First, I utilized the facebook/detr-resnet-101 model for object detection. This model provided me with the locations and labels of the objects in the image, along with confidence scores. The confidence scores for the detected persons ranged from 0.9057888984680176 to 0.997793436050415, indicating a high level of accuracy. After obtaining the detection results, I proceeded to perform image segmentation using the facebook/maskformer-swin-base-coco model on the generated image /images/4516.jpg. The segmentation results can be found in the generated image /images/f2e8.jpg. I chose these models because they are well-suited for object detection and image segmentation tasks, respectively. The object detection model, facebook/detr-resnet-101, is particularly effective at identifying objects within an image, while the facebook/maskformer-swin-base-coco model excels at segmenting objects from the background.}
}
\end{minipage}
\hfill
\begin{minipage}[t]{0.48\textwidth}
\textbf{Explanation B}

\vspace{0.5em}
\fbox{
\parbox[t][0.30\textheight][t]{0.95\linewidth}{To accomplish this task, I used a two-stage approach. First, I utilized the facebook/detr-resnet-101 model for object detection. This model provided me with the locations and labels of the objects in the image, along with confidence scores. The object detection results are as follows: Person 1: score = 0.9308, box = (128, 452, 139, 468), Person 2: score = 0.9978, box = (198, 491, 219, 546), Person 3: score = 0.9807, box = (54, 446, 66, 478), Person 4: score = 0.9057, box = (337, 411, 347, 434), Person 5: score = 0.9673, box = (364, 419, 374, 441). The confidence scores for the detected persons ranged from 0.9057888984680176 to 0.997793436050415, indicating a high level of accuracy. After obtaining the detection results, I proceeded to perform image segmentation using the facebook/maskformer-swin-base-coco model on the generated image /images/4516.jpg. The segmentation results can be found in the generated image /images/f2e8.jpg. I chose these models because they are well-suited for object detection and image segmentation tasks, respectively. The object detection model, facebook/detr-resnet-101, is particularly effective at identifying objects within an image, while the facebook/maskformer-swin-base-coco model excels at segmenting objects from the background. Please note that the accuracy of these results depends on the quality of the input image and the capabilities of the models used.}
}
\end{minipage}

\vspace{1em}

\textbf{Evaluation}
\begin{itemize}
    \item \textbf{Better explanation:} \hspace{1em} $\square$ A \hspace{1em} $\square$ B \hspace{1em} $\square$ Equal
    \item \textbf{Grade A:} \underline{\hspace{2cm}} \hspace{1em} \textbf{Grade B:} \underline{\hspace{2cm}}
    \item \textbf{Reason:} \\[4.5em]
\end{itemize}

\vspace{1em}
\hrule
\vspace{1em}


% item=16 | prompt_id=39 | iteration=1 | compare_mode=base_vs_combined | judge_winner=combined | A_source=Combined | B_source=Base
\subsubsection*{Item 16}
\textbf{Prompt ID:} 39 \hfill \textbf{Iteration:} 1

\textbf{Task Prompt:} Detect pose in 'examples/runner\_pose.jpg' and generate image 'a person jumping'.

\vspace{0.75em}
\noindent
\begin{minipage}[t]{0.48\textwidth}
\textbf{Explanation A}

\vspace{0.5em}
\fbox{
\parbox[t][0.30\textheight][t]{0.95\linewidth}{Given the provided image 'examples/runner\_pose.jpg', I attempted to detect the pose using the openpose-control model, but unfortunately, it was not loaded, and I couldn't proceed with the pose detection. To resolve this issue, ensure the openpose-control model is properly loaded before attempting pose detection. Nevertheless, I was able to generate an image of 'a person jumping' using the runwayml/stable-diffusion-v1-5 model. The generated image is located at '/images/f2e9.jpg'. For further analysis, consider using the generated image as input for the pose detection task or explore other image classification models to gather more relevant information.}
}
\end{minipage}
\hfill
\begin{minipage}[t]{0.48\textwidth}
\textbf{Explanation B}

\vspace{0.5em}
\fbox{
\parbox[t][0.30\textheight][t]{0.95\linewidth}{Given the provided image 'examples/runner\_pose.jpg', I attempted to detect the pose using the openpose-control model, but unfortunately, it was not loaded, and I couldn't proceed with the pose detection. Nevertheless, I was able to generate an image of 'a person jumping' using the runwayml/stable-diffusion-v1-5 model. The generated image is located at '/images/f2e9.jpg'. I used the microsoft/beit-base-patch16-224-pt22k-ft22k model for image classification, which predicted the image to be related to 'mikado, tenno' with a score of 0.19678522646427155, but this result is not directly relevant to the pose detection or image generation tasks.}
}
\end{minipage}

\vspace{1em}

\textbf{Evaluation}
\begin{itemize}
    \item \textbf{Better explanation:} \hspace{1em} $\square$ A \hspace{1em} $\square$ B \hspace{1em} $\square$ Equal
    \item \textbf{Grade A:} \underline{\hspace{2cm}} \hspace{1em} \textbf{Grade B:} \underline{\hspace{2cm}}
    \item \textbf{Reason:} \\[4.5em]
\end{itemize}

\vspace{1em}
\hrule
\vspace{1em}


% item=17 | prompt_id=19 | iteration=7 | compare_mode=base_vs_retrieved | judge_winner=retrieved | A_source=Retrieved | B_source=Base
\subsubsection*{Item 17}
\textbf{Prompt ID:} 19 \hfill \textbf{Iteration:} 7

\textbf{Task Prompt:} Determine the number of cars present in 'examples/traffic\_scene.jpg'.

\vspace{0.75em}
\noindent
\begin{minipage}[t]{0.48\textwidth}
\textbf{Explanation A}

\vspace{0.5em}
\fbox{
\parbox[t][0.30\textheight][t]{0.95\linewidth}{To count the number of cars in the given image, I used the object detection model 'facebook/detr-resnet-101'. This model detects objects in an image and returns their labels, confidence scores, and bounding box coordinates. From the inference results, I filtered out the objects labeled as 'car' and counted them. The model detected 17 cars with confidence scores ranging from 0.9043 to 0.9984. The image with the detected objects is saved at '[PATH]'. Please note that the accuracy of the model depends on various factors, including the quality of the image and the complexity of the scene. If the image is unclear or the scene is too complex, the model may not detect all the cars accurately.}
}
\end{minipage}
\hfill
\begin{minipage}[t]{0.48\textwidth}
\textbf{Explanation B}

\vspace{0.5em}
\fbox{
\parbox[t][0.30\textheight][t]{0.95\linewidth}{To determine the number of cars in the image, I used the object detection model 'facebook/detr-resnet-101'. This model detected multiple objects in the image, including cars, trucks, and people. I filtered the results to only include objects labeled as 'car'. The model detected 17 objects with this label, with confidence scores ranging from 0.9043 to 0.9984. The inference results also included bounding box coordinates for each detected object, which can be used to locate the cars in the image. The generated image with bounding boxes can be found at '/images/1a93.jpg'.}
}
\end{minipage}

\vspace{1em}

\textbf{Evaluation}
\begin{itemize}
    \item \textbf{Better explanation:} \hspace{1em} $\square$ A \hspace{1em} $\square$ B \hspace{1em} $\square$ Equal
    \item \textbf{Grade A:} \underline{\hspace{2cm}} \hspace{1em} \textbf{Grade B:} \underline{\hspace{2cm}}
    \item \textbf{Reason:} \\[4.5em]
\end{itemize}

\vspace{1em}
\hrule
\vspace{1em}


% item=18 | prompt_id=3 | iteration=3 | compare_mode=base_vs_retrieved | judge_winner=base | A_source=Retrieved | B_source=Base
\subsubsection*{Item 18}
\textbf{Prompt ID:} 3 \hfill \textbf{Iteration:} 3

\textbf{Task Prompt:} Generate an image from the text 'a medieval castle under the moonlight'.

\vspace{0.75em}
\noindent
\begin{minipage}[t]{0.48\textwidth}
\textbf{Explanation A}

\vspace{0.5em}
\fbox{
\parbox[t][0.30\textheight][t]{0.95\linewidth}{I used the runwayml/stable-diffusion-v1-5 model to generate the image. This model is a type of generative model that uses diffusion-based image synthesis to create high-quality images from text prompts. The model confidence score is not available for this task. The generated image is located at a local file path and may not be accessible to you. If you need to access the image, I can try to provide a publicly accessible URL or describe the image to you in more detail. To proceed, please let me know if you would like me to share the image or provide more information about it.}
}
\end{minipage}
\hfill
\begin{minipage}[t]{0.48\textwidth}
\textbf{Explanation B}

\vspace{0.5em}
\fbox{
\parbox[t][0.30\textheight][t]{0.95\linewidth}{I used the runwayml/stable-diffusion-v1-5 model to generate the image. This model is a type of generative model that uses diffusion-based image synthesis to create high-quality images from text prompts. The model confidence score for this image is 0.85, indicating a high level of confidence in the generated image. The generated image can be found at file:///images/fa2e.jpg. Please note that this path is a local file path and may not be accessible to you. If you need to access the image, I can try to provide a publicly accessible URL or describe the image in more detail.}
}
\end{minipage}

\vspace{1em}

\textbf{Evaluation}
\begin{itemize}
    \item \textbf{Better explanation:} \hspace{1em} $\square$ A \hspace{1em} $\square$ B \hspace{1em} $\square$ Equal
    \item \textbf{Grade A:} \underline{\hspace{2cm}} \hspace{1em} \textbf{Grade B:} \underline{\hspace{2cm}}
    \item \textbf{Reason:} \\[4.5em]
\end{itemize}

\vspace{1em}
\hrule
\vspace{1em}


% item=19 | prompt_id=69 | iteration=3 | compare_mode=base_vs_retrieved | judge_winner=retrieved | A_source=Retrieved | B_source=Base
\subsubsection*{Item 19}
\textbf{Prompt ID:} 69 \hfill \textbf{Iteration:} 3

\textbf{Task Prompt:} Detect objects and perform segmentation on 'examples/warehouse.jpg'.

\vspace{0.75em}
\noindent
\begin{minipage}[t]{0.48\textwidth}
\textbf{Explanation A}

\vspace{0.5em}
\fbox{
\parbox[t][0.30\textheight][t]{0.95\linewidth}{To accomplish this task, I used a two-stage approach. First, I utilized the facebook/detr-resnet-101 model for object detection. This model provided me with the locations and labels of the objects in the image, along with confidence scores. The object detection results are as follows: Person 1: score = 0.9308, box = (128, 452, 139, 468), Person 2: score = 0.9978, box = (198, 491, 219, 546), Person 3: score = 0.9807, box = (54, 446, 66, 478), Person 4: score = 0.9057, box = (337, 411, 347, 434), Person 5: score = 0.9673, box = (364, 419, 374, 441). The confidence scores for the detected persons ranged from 0.9057888984680176 to 0.997793436050415, indicating a high level of accuracy. After obtaining the detection results, I proceeded to perform image segmentation using the facebook/maskformer-swin-base-coco model on the generated image [PATH] The segmentation results can be found in the generated image [PATH] I chose these models because they are well-suited for object detection and image segmentation tasks, respectively. The object detection model, facebook/detr-resnet-101, is particularly effective at identifying objects within an image, while the facebook/maskformer-swin-base-coco model excels at segmenting objects from the background. Please note that the accuracy of these results depends on the quality of the input image and the capabilities of the models used.}
}
\end{minipage}
\hfill
\begin{minipage}[t]{0.48\textwidth}
\textbf{Explanation B}

\vspace{0.5em}
\fbox{
\parbox[t][0.30\textheight][t]{0.95\linewidth}{To accomplish this task, I used two models in sequence. First, I utilized the 'facebook/detr-resnet-101' model for object detection. This model identified 5 persons in the image with confidence scores ranging from 0.9057 to 0.9978. The detection results were then used to generate a new image, which is saved at 'file:///images/cc27.jpg'. Next, I applied the 'facebook/maskformer-swin-base-coco' model to perform image segmentation on the generated image. The segmented image is saved at 'file:///images/dbd8.jpg'. I chose these models because they are suitable for the tasks of object detection and image segmentation, respectively. The 'facebook/detr-resnet-101' model is particularly effective for detecting objects in images, while the 'facebook/maskformer-swin-base-coco' model is well-suited for segmenting images.}
}
\end{minipage}

\vspace{1em}

\textbf{Evaluation}
\begin{itemize}
    \item \textbf{Better explanation:} \hspace{1em} $\square$ A \hspace{1em} $\square$ B \hspace{1em} $\square$ Equal
    \item \textbf{Grade A:} \underline{\hspace{2cm}} \hspace{1em} \textbf{Grade B:} \underline{\hspace{2cm}}
    \item \textbf{Reason:} \\[4.5em]
\end{itemize}

\vspace{1em}
\hrule
\vspace{1em}


% item=20 | prompt_id=75 | iteration=4 | compare_mode=base_vs_retrieved | judge_winner=base | A_source=Base | B_source=Retrieved
\subsubsection*{Item 20}
\textbf{Prompt ID:} 75 \hfill \textbf{Iteration:} 4

\textbf{Task Prompt:} Combine captions, detections, and segmentation to describe 'examples/beach\_party.jpg'.

\vspace{0.75em}
\noindent
\begin{minipage}[t]{0.48\textwidth}
\textbf{Explanation A}

\vspace{0.5em}
\fbox{
\parbox[t][0.30\textheight][t]{0.95\linewidth}{There are 17 people and 5 boats detected in the image with high confidence scores. Two models were used for this task: facebook/detr-resnet-101 for object detection and facebook/maskformer-swin-base-coco for image segmentation. The object detection model detected people and boats with high confidence scores, while the image segmentation model generated a segmented image of the scene. 
The object detection model detected the following objects: 
- 17 people with confidence scores ranging from 0.9037 to 0.9972 
- 5 boats with confidence scores ranging from 0.9556 to 0.9920 
- 4 umbrellas with confidence scores ranging from 0.9450 to 0.9765 
The image segmentation model generated a segmented image of the scene, which can be found at /images/95dc.jpg. 
The object detection model generated an image with bounding boxes around the detected objects, which can be found at /images/ad6e.jpg.}
}
\end{minipage}
\hfill
\begin{minipage}[t]{0.48\textwidth}
\textbf{Explanation B}

\vspace{0.5em}
\fbox{
\parbox[t][0.30\textheight][t]{0.95\linewidth}{There are 17 people and 5 boats detected in the image with high confidence scores. Two models were used for this task: facebook/detr-resnet-101 for object detection and facebook/maskformer-swin-base-coco for image segmentation. The object detection model detected people and boats with high confidence scores, while the image segmentation model generated a segmented image of the beach party scene. 
The object detection model detected the following objects: 
- 17 people with confidence scores ranging from 0.9037 to 0.9972 
- 5 boats with confidence scores ranging from 0.9556 to 0.9920 
- 4 umbrellas with confidence scores ranging from 0.9450 to 0.9765 
The image segmentation model generated a segmented image of the beach party scene, which can be found at [PATH] 
The object detection results can be found at [PATH]}
}
\end{minipage}

\vspace{1em}

\textbf{Evaluation}
\begin{itemize}
    \item \textbf{Better explanation:} \hspace{1em} $\square$ A \hspace{1em} $\square$ B \hspace{1em} $\square$ Equal
    \item \textbf{Grade A:} \underline{\hspace{2cm}} \hspace{1em} \textbf{Grade B:} \underline{\hspace{2cm}}
    \item \textbf{Reason:} \\[4.5em]
\end{itemize}

\vspace{1em}
\hrule
\vspace{1em}

Prompt 55 iteration 5

\end{document}